% Clase 10 — que dicen los dos potenciales que quedaron.
% CUIDADO CON EL ALCANCE: el docente cierra la clase diciendo que la discusion
% de la solucion (campo electrico, cargas) queda para la clase siguiente. Esta
% figura por lo tanto NO discute lineas de campo ni cargas inducidas: se limita
% a leer la estructura de los dos phi que el si escribio en el pizarron, que es
% un gradiente de distancia. La nota al pie deja constancia de la postergacion.
% Las flechas interiores van deliberadamente mas cortas que las exteriores: esa
% diferencia de longitud ES el contenido (el campo adentro esta reducido por
% 2/(K+1) < 1, porque K > 1).
\begin{tikzpicture}[scale=1]

  % ---------- panel izquierdo: phi_1 = aplicado + correccion ----------
  \begin{scope}[shift={(0,0)}]
    \node[figlbl] at (0,3.0) {afuera: $\phi_1$};
    \fill[figblue!5] (-2.4,-2.4) rectangle (2.4,2.4);
    \fill[figamber!14] (0,0) circle (0.95);
    \draw[figamber, line width=0.9pt] (0,0) circle (0.95);

    \foreach \y in {-1.9,-0.95,0.95,1.9}{
      \draw[figvecaux] (-2.3,\y) -- (-1.3,\y);}
    \node[figlblaux, figgray] at (-1.8,2.25) {$-E_0 r\cos\theta$};

    % el termino correctivo, dipolar en 2D
    \draw[figvec, figred] (1.25,0.62) -- (2.1,0.42);
    \draw[figvec, figred] (1.25,-0.62) -- (2.1,-0.42);
    \draw[figvec, figred] (1.35,0) -- (2.25,0);
    \node[figlblaux, figred, align=center] at (1.15,-1.72)
      {$a^2E_0\dfrac{K-1}{K+1}\dfrac{\cos\theta}{r}$};
    \node[figlblaux, figred] at (1.15,-2.25) {decae como $1/r$};
  \end{scope}

  % ---------- panel derecho: phi_2 uniforme y reducido ----------
  \begin{scope}[shift={(6.2,0)}]
    \node[figlbl] at (0,3.0) {adentro: $\phi_2$};
    \fill[figblue!5] (-2.4,-2.4) rectangle (2.4,2.4);
    \fill[figamber!14] (0,0) circle (0.95);
    \draw[figamber, line width=0.9pt] (0,0) circle (0.95);

    % campo exterior, para comparar longitudes
    \foreach \y in {-1.9,1.9}{
      \draw[figvecaux] (-2.3,\y) -- (-1.0,\y);}
    \node[figlblaux, figgray] at (-1.65,2.25) {$E_0$};

    % campo interior: uniforme y MAS CORTO
    \foreach \y in {-0.45,0,0.45}{
      \draw[figvec, figred, line width=1.0pt] (-0.55,\y) -- (0.25,\y);}
    \node[figlblaux, figred, align=center] at (0.15,-1.72)
      {$\vec E_2=\dfrac{2E_0}{K+1}\,\hat x$};
    \node[figlblaux, figred] at (0.15,-2.25) {uniforme y menor};
  \end{scope}

  \node[figlbl, align=center] at (3.1,-3.3)
    {$\phi_2$ es lineal en $x$: el campo interior es \textbf{uniforme}.};
  \node[figlbl, align=center] at (3.1,-3.8)
    {Y como $K>1$, el factor $2/(K+1)$ es menor que uno:};
  \node[figlbl, align=center] at (3.1,-4.3)
    {el diel\'ectrico \emph{apantalla} el campo aplicado.};

\end{tikzpicture}
